\scanpage{90}
équivalentes

a) Cet hom est injectif

b) La restriction de cet hom à $I_{s'}\simeq T$ est de la forme $\kappa_i\circ(\mu\circ\operatorname{id}_T)$, où $\kappa_i$ est un hom local $T\longrightarrow\pi$, et $\mu\in\mathbb N^{*}$

c) $\lambda\vert I$ n'est pas l'hom trivial

d) la section envisagée \inserted{$E_{\widetilde f}$ de $E_{U_{\widetilde K}}$ sur $E_{\widetilde K}$} ne \deleted{\unsure{se déduit|se déduise}}\inserted{provient} pas d'une section de $E_{U_{\widetilde O}}$ sur $E_{\widetilde O}$

e) La section rationnelle \deleted{régulière} $\widetilde f$ de $\unsure{\widehat U|\widetilde U}$ sur $\widetilde S$ n'est pas régulière (i.e. la sect. rat. $f$ de $U$ sur $S$ n'est pas définie en $s$)

\deleted{f)} \deleted{Si}\inserted{Désignons par} $\widetilde f$ \deleted{désigne} la section de $X_{\widetilde O}$ sur $\widetilde S$ définie par $\widetilde f$ \unsure{loc.|\deleted{les}}, \deleted{$\exists\,\unsure{i\in I(k)}\,\unsure{\sim|\simeq}$ : \unsure{tel que}} On a
\[
\widetilde f(s)\in I=X_s\setminus U_s.
\]
De plus, dans ce cas le $i$ dans b) n'est autre que $\widetilde f(s)$.

Nous plaçant dans ce cas, \uline{\unsure{ns|nous} pouvons donc identifier le groupe de décomposition} $N_{s'}\simeq E_{\widetilde K}$ à un sous-groupe de $\deleted{\unsure{\pi_s}}E_{X_s}$, \inserted{ext. de $E_{k(s)}$ par $\pi$,} qui est d'indice fini $\mu$ dans \inserted{un} $N(L_i)$. Ainsi, notre scindage \inserted{$E_{U_\eta}$ sur $E_K$, ou} de $E_{U_{\widetilde K}}$ \inserted{$E_{X_s}$} sur $E_{\widetilde K}$, qui ne provient pas d'un scindage de
\margin{89}
